\section{Dimension Theory}

\subsection{Affine and finite morphisms}

\bnote

    For any algebraic variety $X$, we will usually write $k[X]$ for $\O_X(X)$.

\enote

\bdefn

    A morphism $f\colon X\to Y$ of algebraic varieties is
    \benum
        \item {\emphcolor affine} if there is an open affine covering $Y=\bigcup_iU_i$ such that $f^{-1}(U_i)\subseteq X$ is affine for all $i$,
        \item {\emphcolor finite} if there is an open affine covering $Y=\bigcup_iU_i$ such that $f^{-1}(U_i)\subseteq X$ is affine, and $k[f^{-1}U_i]$ is a finite $k[U_i]$-algebra for all $i$.
    \eenum

\edefn

Note that $f\colon X\to Y$ induces a morphism $f^*\colon k[U]\to k[f^{-1}U]$ for all open $U\subseteq Y$.
So in the definition of a finite morphism, we mean that $f^*$ is finite for each $i$.

\blemm

    $f\colon X\to Y$ is affine iff $f^{-1}U\subseteq X$ is affine for every open affine $U\subseteq Y$.

    Similarly $f$ is finite iff $f^{-1}U\subseteq X$ is affine and $k[f^{-1}U]$ is a finite $k[U]$-algebra for every open affine $U\subseteq Y$.

\elemm

From this lemma it is clear that the composition of affine/finite morphisms is affine/finite.

\bdefn

    A property $P$ of morphisms is {\emphcolor local on the target} if for any morphism $f\colon X\to Y$ and an open cover $\set{U_i}$ of $Y$,
    $f$ satisfies $P$ iff restricting $f$ to $f^{-1}U_i\to U_i$ satisfies $P$.

\edefn

We then see that being affine/finite is local on the target.

We state a few properties of morphisms of affine varieties which are proven in homework (they can be proven directly or by translating to the case of affine schemes).

\bprop

    Let $f\colon X\to Y$ be a morphism of affine varieties, with induced morphism $f^*\colon k[Y]\to k[X]$.
    \benum
        \item A closed subset $Z\subseteq X$ is irreducible iff $\I(Z)\subseteq k[X]$ is prime, and it is a point iff $\I(Z)$ is a point.
        \item For every $T\subseteq X$, $(f^*)^{-1}(\I_X(T))=\I_Y(f(T))$.
        \item For every closed $S\subseteq Y$, $\I_X(f^{-1}S)=\gen{f^*\I_Y(S)}$.
        \item $f^*$ is injective iff $f$ is {\emphcolor dominant}, i.e.\ $f(X)\subseteq Y$ is dense.
        \item $f^*$ is surjective iff $f$ is a closed embedding, i.e.\ its image is closed and it forms an isomorphism $X\cong f(X)$.
    \eenum

\eprop

\bcoro

    Let $f\colon X\to Y$ be a morphism of affine varieties which induces $f^*\colon k[Y]\to k[X]$.
    \benum
        \item If a prime ideal $\p\leq k[X]$ corresponds to the closed irreducible subset $Z\subseteq X$, then $(f^*)^{-1}\p\subseteq k[Y]$ corresponds to the
        closure $\cl{f(Z)}\subseteq Y$.
        \item If a closed subset $Z\subseteq Y$ corresponds to the radical ideal $I\leq k[Y]$, then the irreducible components $Z_i$ of $f^{-1}Z$ correspond
        to minimal prime ideals $\p_i\leq k[X]$ such that $(f^*)^{-1}\p_i\supseteq I$.
    \eenum

\ecoro

\bproof

    \benum
        \item From the above proposition,
        $$ (f^*)^{-1}\I_X(Z) = \I_Y(f(Z)) = \I_Y(\cl{f(Z)}) , $$
        so it corresponds to $\cl{f(Z)}$.
        \item $Z_i$ is the maximal closed irreducible subset of $X$ such that $f(Z_i)\subseteq Z$, equivalently $\cl{f(Z_i)}\subseteq Z$.
        From (1) this corresponds to $(f^*)^{-1}\p_i$ for prime $\p_i$ containing $\I_Y(Z)=I$.
        \qed
    \eenum

\eproof

\bcoro

    \benum
        \item For an affine variety $X$, the correspondence $x\mapsto\I_X(x)$ is a bijection between points of $X$ and maximal ideals of $k[X]$.
        \item For a morphism $f\colon X\to Y$ inducing $f^*\colon k[Y]\to k[X]$, then for every maximal $\m\leq k[X]$, $(f^*)^{-1}\m\leq k[Y]$
        is maximal.
        Furthermore, under the correspondence of (1), the map $x\mapsto f(x)$ corresponds to $\m\mapsto(f^*)^{-1}\m$.
    \eenum

\ecoro

\bproof

    Immediate.
    \qed

\eproof

\blemm

    For algebraic varieties, the following hold.
    \benum
        \item If $\iota\colon Z\inj X$ is a closed embedding, it is finite.
        \item For every $f\in k[X]$, the inclusion $j\colon D_X(f)\inj X$ is affine.
    \eenum

\elemm

\bproof

    For (1) we may assume $Z\subseteq X$ is a closed subset.
    Then for every affine $U\subseteq X$, $\iota^{-1}U=U\cap Z\subseteq U$ is a closed subspace of an affine variety, thus it itself is affine.
    So $\iota\colon\iota^{-1}U\to U$ is a closed embedding and thus $\iota^*\colon k[U]\to k[\iota^{-1}U]$ is surjective, and in particular
    finite (since a surjective map of rings is finite).

    For (2), for $U\subseteq X$ affine, $j^{-1}U=D_U(f)$ is affine.
    \qed

\eproof

\blemm

    Let $f\colon X\to Y$ be a finite morphism, then for every closed $Z\subseteq X$ and every closed $f(Z)\subseteq W$, $f|_Z\colon Z\to W$
    is finite.

\elemm

\bproof

    If $X,Y$ are affine, then $k[X]$ is a finite $k[Y]$-algebra.
    For a closed $Z\subseteq X$, by the above lemma $k[Z]$ is a finite $k[X]$-algebra and thus a finite $k[Y]$-algebra.
    Since $Z\to Y$ factors as $Z\to W\to Y$, $k[Y]\to k[Z]$ factors as $k[Y]\to k[W]\to k[Z]$, and since $k[Y]\to k[W]$ is finite as well, we see that
    $k[W]\to k[Z]$ must be finite.

    In general, let $Y=\bigcup_iU_i$ be an affine covering such that $f^{-1}U_i\subseteq X$ is affine and $k[f^{-1}U_i]$ is a finite $k[U_i]$-algebra.
    Then $W=\bigcup_iW\cap U_i$ is an open affine covering of $W$, and $f|_Z^{-1}(U_i\cap W)=f^{-1}(U_i)\cap Z$ is closed inside $f^{-1}U_i$, thus affine.
    So from the previous paragraph, $k[f|_Z^{-1}(U_i\cap W)]$ is a finite $k[U_i\cap W]$-algebra, showing that $f|_Z$ is finite.
    \qed

\eproof

We recall the following results, proven in our section on integral and finite algebras.

\blemm

    Let $\phi\colon A\to B$ be an integral (in particular finite) morphism of rings.
    \benum
        \item A prime ideal $\q\leq B$ is maximal iff $\phi^{-1}\q\leq A$ is maximal.
        \item If $\q_1\subsetneq\q_2\leq B$ are prime, then $\phi^{-1}\q_1\subsetneq\phi^{-1}\q_2$.
        \item If $\phi$ is injective, then $\phi^*\colon\spec(B)\to\spec(A)$ is surjective: for every prime $\p\leq A$ there is a prime $\q\leq B$ such that
        $\phi^{-1}\q=\p$.
    \eenum

\elemm

\bprop

    Let $f\colon X\to Y$ be a finite morphism, then
    \benum
        \item $f$ is a closed map, meaning it maps closed subsets of $X$ to closed subsets of $Y$.
        \item if $Z_1\subsetneq Z_2\subseteq X$ are closed irreducible, then $f(Z_1)\subsetneq f(Z_2)$.
    \eenum

\eprop

\bproof

    First we show (1) and (2) in the case that $X,Y$ are affine.

    The induced morphism $f|_Z\colon Z\to\cl{f(Z)}$ is finite as well, from an above proposition, and dominant by definition.
    So replacing $f$ with $f|_Z$, it is sufficient to show that dominant finite morphisms are surjective.

    The map $x\mapsto f(x)$ corresponds to $\m\mapsto(f^*)^{-1}\m$, so we show the second is surjective.
    Let $\m\leq k[Y]$ be maximal, then since $f$ is dominant $f^*$ is injective and finite, and by the above lemma this means there is a prime $\p\leq k[X]$
    such that $(f^*)^{-1}\p=\m$.
    And by the same lemma, this means $\p$ is maximal, so $(f^*)^{-1}$ is indeed surjective.

    Since $Z_1\subsetneq Z_2$ are closed irreducible, $\I_X(Z_1)\supsetneq\I_X(Z_2)$ are prime (the inequality is due to $\I_X$ being a bijection on closed subsets).
    Since $k[X]$ is a finite $k[Y]$-algebra, by the above lemma
    $$ \I_Y(f(Z_1)) = (f^*)^{-1}\I_X(Z_1) \supsetneq (f^*)^{-1}\I_X(Z_2) = \I_Y(f(Z_2)) , $$
    and thus $f(Z_1)\subsetneq f(Z_2)$.

    Now in general take $Y=\bigcup_iU_i$ an affine covering such that $f^{-1}U_i\subseteq X$ is affine and $k[f^{-1}U_i]$ is a finite $k[U_i]$-algebra.
    Since $Z\subseteq X$ is closed, $Z\cap f^{-1}U_i\subseteq f^{-1}U_i$ is closed.
    Then as in the affine case, $f(Z\cap f^{-1}U_i)\subseteq f(U_i)$ is closed.
    Since $f(Z\cap f^{-1}U_i)=f(Z)\cap U_i$, we see that $f(Z)$ is closed as desired.

    This is because being closed is ``local'': if $Z\cap U_i$ is closed for all $i$ in an open covering then $Z$ is closed.
    And this is because if we let $W=Y-Z$, we see $W\cap U_i$ is open for all $i$ and thus $W$ is open and so $Z$ is closed.

    For (2), let $i$ be such that $f(Z_2)\cap U_i\neq\emptyset$.
    Then $Z_1\cap f^{-1}U_i\subseteq Z_2\cap f^{-1}U_i$ are closed in $f^{-1}U_i$.
    And $Z_2\cap f^{-1}U_i\subseteq Z_2$ is open and nonempty, thus dense.
    Since $Z_1\cap f^{-1}U_i\subseteq Z_1\subsetneq Z_2$, it is not dense in $Z_2$, so $Z_1\cap f^{-1}U_i\subsetneq Z_2\cap f^{-1}U_i$.
    These are closed irreducible subsets of an affine variety, so $f(Z_1\cap f^{-1}U_i)\subsetneq f(Z_2\cap f^{-1}U_i)$,
    and thus we must have $f(Z_1)\subsetneq f(Z_2)$.
    \qed

\eproof

\bcoro

    If $X$ is an algebraic variety, its points are all closed.

\ecoro

\bproof

    For $x\in X$, $x$ is a point of an affine variety and thus itself an affine variety.
    Its ideal is maximal, which means $k[\set x]=k$ and so $\set x\inj X$ is finite as $k$ is a finite $k[X]$-algebra.
    So by the above proposition the inclusion is a closed map, and thus $x$ is closed in $X$.

    Alternatively if we choose an affine cover $U_i$ of $X$, then $\set x\cap U_i$ is either empty or $\set x\subseteq U_i$.
    In each case it is closed, and thus $\set x$ is closed.
    \qed

\eproof

\bdefn

    A morphism of varieties $f\colon X\to Y$ is {\emphcolor quasi-finite} if it has finite fibers: $f^{-1}y$ is a finite set for all $y\in Y$.

\edefn

\bcoro

    Every finite morphism of varieties is quasi-finite.

\ecoro

\bproof

    Let $f\colon X\to Y$ be finite, and $y\in Y$.
    It is sufficient to show that every irreducible component of $f^{-1}y$ is a point, as varieties have finitely many irreducible components.
    Let $Z$ be an irreducible component, and let $x\in Z$, then $\set x\subseteq Z$ are closed irreducible sets and $f(Z)=f(\set x)=y$.
    By the proposition, we must have $\set x=Z$.
    \qed

\eproof

\bcoro

    \benum
        \item For every affine variety $X$ there is a finite surjective morphism $X\to\bA^n$ such that $k[X]$ is a finite $k[\bA^n]$-algebra.
        \item For every non-constant function $f\in k[\bA^n]-k$, there is a finite surjective morphism $Z(f)\to\bA^{n-1}$ such that $k[Z(f)]$ is a finite
        $k[\bA^{n-1}]$-algebra.
    \eenum

\ecoro

\bproof

    By the Noether normalization lemma, since $k[X]$ is fg there is a finite embedding $k[\bA^n]\inj k[X]$ for some $n$.
    Thus $X\to\bA^n$ is finite and dominant.
    It is then closed, and since it is dominant it is surjective.

    For (2), we have $k[Z(f)]=k[\bA^n]/(f)$.
    Then in Noether normalization our injective embedding was $k[\bA^{n-1}]\to k[bA^n]/(f)$.
    \qed

\eproof

\subsection{Dimension theory}

\bdefn

    Let $X$ be a topological space.
    A {\emphcolor chain of irreducible subsets} is a sequence $\emptyset\neq Z_0\subsetneq Z_1\subsetneq\cdots\subsetneq Z_n\subseteq X$
    where each $Z_i$ is irreducible.
    The length of such a chain is $n$.

    Then we define the {\emphcolor (Krull) dimension} of $X$ to be the supremum of the lengths of all chains of irreducible subsets.
    It is denoted $\dim X$.
    For $X=\emptyset$ we set $\dim X=-1$.

\edefn

\bnote

    For Hausdorff $X$, the only irreducible subsets are the points, and so $\dim X=0$.
    We will thus consider only Noetherian topological spaces.

    The Krull dimension of a Noetherian topological space may be infinite.
    While each chain is finite (and can only be extended finitely), there need not be a bound on the length of all chains.
    We will nevertheless show that for algebraic varieties the dimension is finite.

\enote

\blemm

    If $Y$ is irreducible and $X\subseteq Y$ is closed with $\dim X=\dim Y<\infty$, then $X=Y$.

\elemm

\bproof

    We show that if $\dim X=n$ and $X\subsetneq Y$ then $\dim Y\geq n+1$.
    By definition, there is a chain of irreducible subsets of $X$ of length $n$: $\emptyset\neq X_0\subsetneq\cdots\subsetneq X_n$.
    These are then closed irreducible in $Y$, and we can append $Y$ to this chain to obtain a chain of length $n+1$.
    \qed

\eproof

\blemm

    Let $X$ be an arbitrary topological space.
    \benum
        \item If $Y\subseteq X$ is a subspace, $\dim Y\leq\dim X$.
        \item If $X=\bigcup_iX_i$ is a finite union of closed subsets, then $\dim X=\max_i\dim X_i$.
        \item If $X=\bigcup_iU_i$ is an open covering, then $\dim X=\sup_i\dim U_i$.
    \eenum

\elemm

\bproof

    For (1), we note that if $Y_n\subseteq Y$ is closed irreducible then $Y_n=Y\cap X_n$ for $X_n\subseteq X$ irreducible.
    Indeed, $Y_n=Y\cap X_n$ for $X_n$ closed.
    Then $X_n$ is the union of its irreducible components, and then $Y_n$ is the union of the intersection of $Y$ with these irreducible components.
    There are all closed in $Y_n$, and thus this union cannot be proper, so $Y_n$ is equal to the intersection of $Y$ with one of these irreducible components.

    So every chain in $Y$ gives rise to a chain in $X$, showing $\dim Y\leq\dim X$.

    For (2) and (3), by (1) we have that $\max_i\dim X_i,\sup_i\dim U_i\leq\dim X$.
    Conversely let $\emptyset\neq Z_0\subsetneq\cdots\subsetneq Z_n$ be a chain of irreducible subsets of $X$.
    Then for each $j$,
    $$ Z_j = \bigcup_iZ_j\cap X_i . $$
    Since $Z_j$ is irreducible and this is a finite union of closed subsets, this means there is an $i$ where $Z_j=Z_j\cap X_i$ and so $Z_j\subseteq X_i$.
    For $j=n$ this shows that this chain is contained in some $X_i$ so $\dim X\leq\max_i\dim X_i$.

    For (3) we note that open subsets of irreducible subsets are irreducible (since irreducibility is equivalent to open subsets having nonempty intersection).
    Since $Z_0$ is nonempty there is some $U_i$ where $Z_0\cap U_i\subseteq U_i$ is nonempty.
    In such a case we also have $Z_j\cap U_i\subsetneq Z_{j+1}\cap U_i$ since $Z_{j+1}\cap U_i$ is dense in $Z_{j+1}$ while $Z_j$ is not.
    So we get an irreducible chain $\emptyset\neq Z_0\cap U_i\subsetneq\cdots\subsetneq Z_n\cap U_i$ of $U_i$, so $\sup_i\dim U_i\geq\dim X$.
    \qed

\eproof

\blemm

    If $f\colon X\to Y$ is a finite surjective morphism of algebraic varieties.
    Then $\dim X=\dim Y$.

\elemm

\bproof

    If $\emptyset\neq X_0\subsetneq\cdots\subsetneq X_n$ is a chain of closed irreducible subsets of $X$, then by our previous results
    $\emptyset\neq f(X_0)\subsetneq\cdots\subsetneq f(X_n)$ is a chain of closed irreducible subsets of $Y$, so $\dim Y\geq\dim X$.

    Conversely let $\emptyset\neq Y_0\subsetneq\cdots\subsetneq Y_n$ be a chain of closed irreducible subsets of $Y$.
    We claim that there is a chain $\emptyset\neq X_0\subsetneq\cdots\subsetneq X_n$ of closed irreducible subsets of $X$ such that $f(X_i)=Y_i$ for all $i$.

    Write $f^{-1}Y_n$ as a union of its irreducible components, $f^{-1}Y_n=\bigcup_jK_j$.
    Since $f$ is surjective, $Y_n=\bigcup f(K_j)$ and $f$ is a closed map so $f(K_j)\subseteq Y$ is closed.
    Thus $f(K_j)=Y_n$ for some $j$; let $X_n=K_j$.

    The restriction $f|_{X_n}\colon X_n\to Y_n$ is finite surjective, so inducting on $n$ we construct $\emptyset\neq X_0\subsetneq\cdots\subsetneq X_n$
    as desired.
    \qed

\eproof

\bnote

    Later we will be able to show that for a morphism of varieties $f\colon X\to Y$,
    \benum
        \item if $f$ is dominant, then $\dim X\geq\dim Y$,
        \item if $f$ is quasi-finite, then $\dim X\leq\dim Y$.
    \eenum
    So if $f$ is dominant and quasi-finite (e.g.\ surjective and finite) then $\dim X=\dim Y$.

\enote

\bthrm

    \benum
        \item For $n\geq0$, $\dim\bA^n=n$.
        \item For $n>0$ and $f\in k[\bA^n]-k$, $\dim Z(f)=n-1$.
    \eenum

\ethrm

\bproof

    We prove both points by induction on $n$.
    For $n=0$ point (2) holds vaccuously and point (1) holds trivially since $\bA^0$ is a single point.
    So assume $n>0$.

    For point (2), there is a finite surjective morphism $Z(f)\to\bA^{n-1}$ so by the above lemma, $\dim Z(f)=\dim\bA^{n-1}$ which is $n-1$ by induction.
    For (1), consider the chain $X_0\subsetneq X_1\subsetneq\cdots\subsetneq X_n$ of $\bA^n$ where
    $$ X_i = Z(x_{i+1},\dots,x_n) . $$
    These are all irreducible as $(x_{i+1},\dots,x_n)$ is prime and the strict inclusion is clear.
    So $\dim\bA^n\geq n$.

    Conversely let $X_0\subsetneq\cdots\subsetneq X_m$ be a chain of closed irreducible subsets of $\bA^n$.
    Now we have that $\dim X_{m-1}\geq m-1$, since this chain gives a chain of length $m-1$ of $X_{m-1}$, and $X_{m-1}\subsetneq\bA^n$.
    Thus there is a $f\in k[\bA^n]-0$ such that $X_{m-1}\subseteq Z(f)$ (because $\I(X_{m-1})\neq0$), and thus we have $f\notin k$.
    So inductively $\dim Z(f)=n-1$ and thus
    $$ m - 1 \leq \dim X_{m-1} \leq \dim Z(f) = n - 1 \implies m\leq n $$
    as desired.
    \qed

\eproof

\bcoro

    The dimension of any algebraic variety is finite.

\ecoro

\bproof

    The dimension of any affine variety is finite, as a subspace of $\bA^n$.
    If $X$ is an algebraic variety, it has a finite cover by open affine subvarieties.
    Thus $X$'s dimension is the maximum of these dimensions, which is finite.
    \qed

\eproof

\bdefn

    Let $L/K$ be a field extension.
    \benum
        \item A subset $S\subseteq L$ is {\emphcolor algebraically independent} over $K$ if for every $s_1,\dots,s_n\in S$ and $0\neq f\in K[x_1,\dots,x_n]$
        we have $f(s_1,\dots,s_n)\neq0$.
        \item A subset $S\subseteq L$ is a {\emphcolor transcendence basis} of $L/K$ if it is algebraically independent over $K$ and $L$ is an algebraic extension
        of $K(S)$.
        \item It can be shown that every field extension has a transcendence basis, and they all have the same cardinality.
        The cardinality of a transcendence basis is called the {\emphcolor transcendence degree} of the extension, and is denoted $\trdeg(L/K)$ or $\trdeg_K(L)$.
    \eenum

\edefn

\bnote

    To prove the existence of a transcendence basis, and the uniqueness of its cardinality, we can use the concept of a {\emphcolor pregeometry}.
    A {\emphcolor pregeometry} is a pair $(S,\cll)$ where $\cll\colon\P(S)\to\P(S)$ is a {\emphcolor closure operation} satisfying the following axioms:
    \benum
        \item {\emphcolor monotonicity}: $\cll$ is monotone increasing and dominates: if $A\subseteq B$ then $A\subseteq\cll(A)\subseteq\cll(B)$.
        \item {\emphcolor idempotent}: $\cll$ is idempotent: $\cll(\cll(A))=\cll(A)$.
        \item {\emphcolor finite character}: for each $a\in\cll(A)$ there is some finite $F\subseteq A$ with $a\in\cll(F)$.
        \item {\emphcolor exchange}: if $b\in\cll(A\cup\set c)-\cll(A)$, then $c\in\cll(A\cup\set b)$.
    \eenum

    For a pregeometry $(S,\cll)$, a subset $D\subseteq S$ is {\emphcolor independent} if $d\notin\cll(D-d)$ for all $d\in D$.
    A {\emphcolor basis} for $S$ is an independent set $D$ such that $\cll(D)=S$.
    Equivalently this is a maximally independent set; and thus it can be shown easily via Zorn that a basis must exist.

    Then we can show that if $A$ is independent and $B$ generates $S$ (meaning $\cll(B)=S$), then $\abs A\leq\abs B$.
    Clearly this implies two bases have the same cardinality.
    We omit the details of the proofs here, they can be found in most standard texts on model theory for example (or my final assignment in model theory).

    Pregeometries are also called {\emphcolor matroids} outside of model theory.
    They generalize the concept of independence found in linear algebra and are used to prove important results in model theory.
    Our main concern here is their use for algebraic independence.

    Given a field extension $L/K$, we define a pregeometry on $L$ by
    $$ \cll(S) = \set{a\in L}[\hbox{$a$ is algebraic over $K(S)$}] . $$
    This is clearly monotonic and dominates, idempotence is easily seen as well.
    Finite character is easy to see as well: if $a$ is algebraic over $S$ then it is the root of some polynomial in $K(S)[x]$.
    Any such polynomial must use finitely many elements from $S$.
    If $a$ is algebraic over $K(S\cup b)$ but not $K(S)$, then there is a polynomial in $K(S\cup b)[x]$ with root $a$, and it must have some coefficient with $b$.
    Then viewing $b$ as the variable and not $a$, we see $b\in\cll(A\cup a)$, so this satisfies exchange.

    If $S$ generates $L$ then every $a\in L=\cll(S)$ is algebraic over $K(S)$, so $L/K(S)$ is algebraic.
    And if $S$ is independent, then let $s_1,\dots,s_n\in S$ and suppose $f(s_1,\dots,s_n)=0$ for $0\neq f\in K[x_1,\dots,x_n]$.
    Take $n$ minimal, so each $x_i$ has nonzero coefficient.
    This shows that $s_1$ is algebraic over $K(s_2,\dots,s_n)$, which contradicts independence.
    So if $S$ is independent in the pregeometry, then it is algebraically independent over $K$ (the converse is readily shown as well).

    Thus a basis of the pregeometry is the same as a transcendence basis of the field extension, and we have proven the result.

\enote

If we have field extensions $M/L/K$, then let $S_0$ be a transcendence basis of $M/L$, and $S_1$ a transcendence basis of $L/K$.
These are disjoint as $S_0\subseteq M-L$, and their union forms a transcendence basis of $M/K$.
So
$$ \trdeg(M/K) = \trdeg(M/L) + \trdeg(L/K) . $$

If $L/K$ is finitely generated, so $L=K(y_1,\dots,y_n)$, then any maximal algebraically independent subset of $y_1,\dots,y_n$ over $K$ is a transcendence basis.
Thus $\trdeg(L/K)\leq n$.

\bdefn

    Let $X$ be an algebraic variety.
    We define $k(X)$ to be the colimit of sections of dense open subsets:
    $$ k(X) = \colim_{U\subseteq X\hbox{ dense}}\O_X(U) . $$
    Explicitly, we can write $k(X)$ in terms of equivalence classes.
    We identify two pairs $(U,f)$ and $(V,g)$ where $f\in\O_X(U),g\in\O_X(V)$ if there is a dense open subset $W\subseteq U\cap V$ such that $f|_W=g|_W$.
    Then $k(X)$ is the quotient
    $$ k(X) = \coprod_{U\subseteq X\hbox{ dense}}\O_X(U)/{\sim} . $$

    A {\emphcolor rational function} is an element of $k(X)$, so an equivalence class.

\edefn

\blemm

    \benum
        \item $k(X)$ has a natural $k$-algebra structure.
        \item For any dense open subset $U\subseteq X$, restriction induces a $k$-algebra isomorphism $k(X)\cong k(U)$.
        \item If $X=\bigcup_{i=1}^nZ_i$ is the decomposition of $X$ into irreducible components, then $k(X)\cong\prod_ik(Z_i)$ where the components of the map
        are restrictions.
        \item If $X$ is irreducible, $k(X)$ is a field.
        \item If $X$ is affine and irreducible, then $k(X)$ is the fraction field of $k[X]$.
        More precisely the map $k[X]\to k(X)$ given by $f\mapsto(X,f)$ factors through an isomorphism $\Frac(k[X])\cong k(X)$.
    \eenum

\elemm

\bproof

    Point (1) is clear as a colimit of $k$-algebras is a $k$-algebra.
    Point (2) is clear since dense open subsets of $U$ are dense in $X$, and form a cofinal subcategory of dense open subsets of $X$.
    This is because $U$ is dense, so for every $W\subseteq X$ open dense we have $W\cap U\subseteq U$ is open dense.

    For (3), define $U_i=Z_i-\bigcup_{j\neq i}Z_j$.
    Since $Z_i$ are irreducible, $U_i\neq\emptyset$, and since $U_i=X-\bigcup_{j\neq i}Z_j$, $U_i$ is open.
    And the $U_i$ are all disjoint.
    Define $U=\bigcup_iU_i$, then we claim $U$ is dense in $X$.

    Indeed, let $V\subseteq X$ be nonempty open.
    Then $V\cap Z_i\neq\emptyset$ for some $i$, so $V\cap Z_i$ is nonempty open in $Z_i$ and thus dense.
    So $V\cap Z_i\cap U_i\neq\emptyset$, and so $V\cap U\neq\emptyset$ meaning $U$ is indeed dense.

    But $U=\bigsqcup_iU_i$ is a disjoint union of irreducible $U_i$, so a dense open subset of $U$ is the same as a disjoint union $W=\bigsqcup_iW_i$
    for $W_i\subseteq U_i$ dense open.
    By the sheaf condition, $k[W]=\prod_ik[W_i]$.
    This says that the functor $W\mapsto\O_X(W)$ is the same as a the product functor $(W_i)\mapsto\prod_ik[W_i]$.
    Colimits and limits commute, so we see
    $$ k(U) \cong \prod_ik(U_i) . $$
    Now, $U\subseteq X$ is dense and $U_i\subseteq Z_i$ is dense, so by (2) this forms an isomorphism
    $$ k(X) \cong \prod_ik(Z_i) . $$

    For (4), note that all nonempty open $U\subseteq X$ are dense.
    For nonzero $f\in k[U]$, $D_U(f)\subseteq X$ is nonempty open and thus dense.
    So $(D_U(f),1/f)\in k(X)$ is a multiplicative inverse for $(U,f)$.

    For (5), since $k[X]\subseteq k(X)$ and $k(X)$ is a field, it remains to show that every element of $k(X)$ can be written as $g/h$ for $g/h\in k[X]$.
    Indeed, the basic open subsets $D_X(h)$ form a basis of $X$, so every $f\in k(X)$ has representative in $k[D_X(h)]=k[X]_h$.
    Such an $f$ is of the form $g/h^k$ for $g\in k[X]$ as desired.
    \qed

\eproof

\bcoro

    If $X$ is a nonempty irreducible affine variety, $\dim X=\trdeg_k(k(X))$.

\ecoro

\bproof

    There is a finite surjective morphism $X\to\bA^n$, so $\dim X=\dim\bA^n=n$.
    Since $k[\bA^n]\subseteq k[X]$ and $k[X]$ is a finite $k[\bA^n]$-algebra, we have that $k(X)$ is a finite extension of $k(\bA^n)=k(x_1,\dots,x_n)$
    (as it is the fraction field of $k[x_1,\dots,x_n]$).
    Since $x_1,\dots,x_n$ are algebraically independent over $k$ and $k(X)/k(x_1,\dots,x_n)$ is finite and thus algebraic, we see that $\trdeg_k(k(X))=n$.
    \qed

\eproof

\bcoro

    If $X$ is an algebraic variety and $U\subseteq X$ is open dense, then $\dim U=\dim X$.

\ecoro

\bproof

    If $X$ is affine and irreducible, then there is a $f\in k[X]$ such that $D_X(f)\subseteq U\subseteq X$.
    Thus $\dim D_X(f)\leq\dim U\leq\dim X$.
    On the other hand $k[D_X(f)] = k[X]_f$, and since $X$ is irreducible so is $D_X(f)$ and thus $k(D_X(f))$ is the fraction field of $k[X]_f$,
    which is just $k(X)$.
    Thus $\dim D_X(f)=\dim X=\trdeg_k(k(X))$ so $\dim U=\dim X$.

    Now suppose $X$ is just affine, and let $X_1,\dots,X_n$ be the irreducible components of $X$.
    Since $U\subseteq X$ is open dense, $U\cap X_i\subseteq X_i$ is open dense and $X_i$ is affine so $\dim(U\cap X_i)=\dim X_i$.
    Then $U$ is the finite union of the closed $U\cap X_i$, so $\dim U=\max_i\dim X_i=\dim X$ (the last equality is because $X_i$ form a finite closed cover of $X$).

    In the general case, choose an open affine covering $X=\bigcup_iX_i$, then $U\cap X_i\subseteq X_i$ is open dense so $\dim(U\cap X_i)=\dim X_i$.
    So
    $$ \dim U = \max_i\dim(U\cap X_i) = \max_i\dim X_i = \dim X , $$
    as desired.
    \qed

\eproof

\bcoro

    If $X$ is a nonempty irreducible variety, then $\dim X=\trdeg_k(k(X))$.

\ecoro

\bproof

    Choose an open affine neighborhood $U\subseteq X$.
    Since $X$ is irreducible, $U$ is open dense.
    Then $k(X)\cong k(U)$ and $\dim U=\dim X$.
    Since $\dim U=\trdeg_k(k(U))$ we are finished.
    \qed

\eproof

If we define $\trdeg_k(0)=-1$ for the zero ring, then the previous corollary holds for the empty variety.

